\documentclass[11pt,oneside]{article}
\newcommand{\TopicTitle}{Apuntes del Tema 11: Attention Mechanisms and Transformers}
\usepackage[a4paper,margin=2.6cm,headheight=15pt]{geometry}
\usepackage{fontspec}
\usepackage{polyglossia}
\setdefaultlanguage{spanish}
\setotherlanguage{english}
\setmainfont{Latin Modern Roman}
\setsansfont{Latin Modern Sans}
\setmonofont{Latin Modern Mono}[Scale=MatchLowercase]
\usepackage{amsmath,amssymb,mathtools}
\usepackage{graphicx}
\usepackage{xcolor}
\usepackage{booktabs}
\usepackage{tabularx}
\usepackage{array}
\usepackage{longtable}
\usepackage{enumitem}
\usepackage{listings}
\usepackage{fancyhdr}
\usepackage{titlesec}
\usepackage{microtype}
\usepackage{hyperref}
\usepackage{bookmark}
\usepackage{caption}

\definecolor{ink}{HTML}{243447}
\definecolor{accent}{HTML}{176B87}
\definecolor{codebg}{HTML}{F3F5F7}
\definecolor{codegray}{HTML}{65727E}

\hypersetup{
  colorlinks=true,
  linkcolor=accent,
  urlcolor=accent,
  citecolor=accent,
  pdftitle={\TopicTitle},
  pdfauthor={Foundations of Machine Learning and Deep Learning}
}
\setcounter{secnumdepth}{3}
\setcounter{tocdepth}{3}
\setlist[itemize]{leftmargin=1.7em,itemsep=2pt,topsep=4pt}
\setlist[enumerate]{leftmargin=1.9em,itemsep=2pt,topsep=4pt}
\setlength{\parindent}{0pt}
\setlength{\parskip}{5pt plus 1pt minus 1pt}
\titleformat{\section}{\Large\sffamily\bfseries\color{ink}}{\thesection}{0.8em}{}
\titleformat{\subsection}{\large\sffamily\bfseries\color{accent}}{\thesubsection}{0.8em}{}
\titleformat{\subsubsection}{\normalsize\sffamily\bfseries\color{ink}}{\thesubsubsection}{0.8em}{}
\titlespacing*{\section}{0pt}{2.0em}{0.7em}
\titlespacing*{\subsection}{0pt}{1.4em}{0.4em}
\titlespacing*{\subsubsection}{0pt}{1.0em}{0.3em}
\lstset{
  backgroundcolor=\color{codebg},
  basicstyle=\ttfamily\small,
  breaklines=true,
  breakatwhitespace=false,
  columns=fullflexible,
  frame=single,
  framerule=0.3pt,
  rulecolor=\color{codegray},
  numbers=left,
  numbersep=8pt,
  numberstyle=\scriptsize\color{codegray},
  showstringspaces=false,
  keepspaces=true,
  tabsize=4,
  keywordstyle=\color{accent}\bfseries,
  commentstyle=\color{codegray}\itshape,
  stringstyle=\color{purple}
}
\captionsetup{font=small,labelfont=bf}
\pagestyle{fancy}
\fancyhf{}
\fancyhead[L]{\small\sffamily\color{codegray}Foundations of Machine Learning and Deep Learning}
\fancyhead[R]{\small\sffamily\color{codegray}\nouppercase{\rightmark}}
\fancyfoot[C]{\thepage}
\renewcommand{\headrulewidth}{0.4pt}
\renewcommand{\footrulewidth}{0pt}
\newcommand{\topiccover}[1]{%
  \begin{titlepage}
    \centering
    \vspace*{2.5cm}
    {\Large\sffamily\color{accent}APUNTES UNIVERSITARIOS\par}
    \vspace{1.5cm}
    {\Huge\sffamily\bfseries\color{ink}Foundations of Machine Learning\\[0.25em]
    and Deep Learning\par}
    \vspace{1.8cm}
    {\rule{0.72\textwidth}{1pt}\par}
    \vspace{1.2cm}
    {\LARGE\sffamily\bfseries #1\par}
    \vfill
    {\large\sffamily\color{codegray}\today\par}
  \end{titlepage}
}

\begin{document}
\topiccover{Apuntes del Tema 11: Attention Mechanisms and Transformers}
\tableofcontents
\clearpage
\markboth{Apuntes del Tema 11: Attention Mechanisms and Transformers}{Apuntes del Tema 11: Attention Mechanisms and Transformers}
\section{Apuntes del Tema 11: Attention Mechanisms and Transformers}

\subsection{1. Introducción y Limitaciones de los Modelos Secuenciales (Seq2Seq)}

\subsubsection{1.1 El Enfoque Seq2Seq Tradicional (RNNs)}

\begin{itemize}

\item \textbf{Modelado Secuencial:} Las redes neuronales recurrentes (RNNs) procesan las secuencias token a token de forma secuencial.

\end{itemize}

\begin{itemize}

\item \textbf{Arquitectura Encoder-Decoder:} En tareas como la traducción automática (p. ej., de inglés a español), el \emph{Encoder} procesa la frase de entrada y condensa su información en un vector de estado oculto final \(h_{T_x}\). Luego, el \emph{Decoder} genera la secuencia traducida token a token condicionado en este vector.

\end{itemize}

\begin{lstlisting}

Entrada (Inglés) ──► [ ENCODER (RNN) ] ──► Vector de Estado H_Tx ──► [ DECODER (RNN) ] ──► Salida (Español)



\end{lstlisting}

\subsubsection{1.2 El Cuello de Botella del Vector de Contexto (\emph{Context Bottleneck})}

\begin{itemize}

\item \textbf{Limitación Principal:} Comprimir toda una frase de longitud variable en un único vector de tamaño fijo \(h_{T_x}\) genera una gran pérdida de información.

\end{itemize}

\begin{itemize}

\item \textbf{Fallo en Secuencias Largas:} Aunque funciona bien para frases cortas (p. ej., \emph{"The black cat"} \(\rightarrow\) \emph{"El gato negro"}), el modelo falla en oraciones largas y complejas.

\end{itemize}

\begin{itemize}

\item \textbf{Dependencias a Largo Plazo:} Las palabras iniciales se "olvidan" a medida que la secuencia crece. Además, no gestiona eficientemente el reordenamiento de palabras entre idiomas.

\end{itemize}

\subsection{2. Mecanismos de Atención en RNNs (Bahdanau et al., 2014)}

\subsubsection{2.1 Filosofía de la Atención}

En lugar de depender de un único vector fijo, la atención permite al \emph{Decoder} acceder a \textbf{todos} los estados ocultos del \emph{Encoder} \(\{h_1, h_2, \dots, h_{T_x}\}\) en cada paso de decodificación.

\begin{itemize}

\item \textbf{Adaptabilidad:} En cada paso \(t\), el \emph{Decoder} calcula qué partes de la frase de entrada son más relevantes para predecir el siguiente token.

\end{itemize}

\begin{itemize}

\item \textbf{Vector de Contexto Dinámico (\(c_t\)):} Se calcula como una suma ponderada de todos los estados del \emph{Encoder}:

\end{itemize}

\[c_t = \sum_{i=1}^{T_x} \alpha_{t,i} h_i\]

\subsubsection{2.2 Formulación Matemática de la Atención de Bahdanau}

Para cada paso de tiempo del decodificador \(i\) con estado anterior \(s_{i-1}\):

\begin{enumerate}

\item **Puntuación de Alineamiento (\emph{Score}):**

\end{enumerate}

\[e_{ij} = a(s_{i-1}, h_j)\]

\emph{(donde \(a\) es una red feedforward)}.

\begin{enumerate}

\item \textbf{Pesos de Atención (Softmax):}

\end{enumerate}

\[\alpha_{ij} = \frac{\exp(e_{ij})}{\sum_{k=1}^{T_x} \exp(e_{ik})}\]

\begin{enumerate}

\item \textbf{Vector de Contexto:}

\end{enumerate}

\[c_i = \sum_{j=1}^{T_x} \alpha_{ij} h_j\]

\begin{enumerate}

\item \textbf{Actualización del Estado y Predicción:}

\end{enumerate}

\[s_i = f(s_{i-1}, y_{i-1}, c_i)\]

\[P(y_i \mid y_{1:i-1}, x) = g(y_{i-1}, s_i, c_i)\]

\subsubsection{2.3 Inconvenientes de la Atención sobre RNNs}

\begin{itemize}

\item \textbf{Dependencia Recurrente:} La computación sigue siendo secuencial (paso a paso), lo que impide la paralelizabilidad masiva en GPUs durante el entrenamiento.

\end{itemize}

\begin{itemize}

\item \textbf{Poco Directo:} La atención solo conecta el \emph{Decoder} con el \emph{Encoder}, pero la información dentro de cada bloque sigue fluyendo de forma secuencial implícita.

\end{itemize}

\subsection{3. Autoatención (\emph{Self-Attention}) en Transformers}

Propuesto por Vaswani et al. (2017) en \emph{"Attention is All You Need"}, elimina por completo la recurrencia. Cada palabra puede atender directamente a cualquier otra palabra de la secuencia sin importar la distancia.

\subsubsection{3.1 Queries, Keys y Values (\(Q, K, V\))}

Para cada token con embedding \(x \in \mathbb{R}^d\), se calculan tres proyecciones mediante matrices de pesos entrenables:

\begin{enumerate}

\item \textbf{Query (\(q = W_Q x \in \mathbb{R}^{d_k}\)):} Lo que el token busca saber sobre otros.

\end{enumerate}

\begin{enumerate}

\item \textbf{Key (\(k = W_K x \in \mathbb{R}^{d_k}\)):} Cómo se identifica o etiqueta el token ante los demás.

\end{enumerate}

\begin{enumerate}

\item \textbf{Value (\(v = W_V x \in \mathbb{R}^{d_v}\)):} La información que el token lleva para compartir.

\end{enumerate}

\subsubsection{3.2 Formulación Matricial: Scaled Dot-Product Attention}

Dada la matriz de entrada \(X \in \mathbb{R}^{T_x \times d_{\text{model}}}\), se proyectan las matrices \(Q = X W_Q\), \(K = X W_K\) y \(V = X W_V\).

\[\text{Attention}(Q, K, V) = \text{Softmax}\left( \frac{Q K^\top}{\sqrt{d_k}} \right) V\]

\begin{itemize}

\item \textbf{Matriz de Similitud (\(S = \frac{Q K^\top}{\sqrt{d_k}}\)):} Mide la afinidad entre las consultas y las claves.

\end{itemize}

\begin{itemize}

\item \textbf{Factor de Escala (\(\sqrt{d_k}\)):} Previene que el producto escalar alcance magnitudes muy grandes que saturarían la función Softmax (provocando gradientes nulos).

\end{itemize}

\begin{lstlisting}

Entrada X (T x d_model)

    ├──► Q = X·W_Q (T x d_k) ──┐

    ├──► K = X·W_K (T x d_k) ──┴──► [ Q·Kᵀ / √d_k ] ──► [ Softmax ] ──┐

    └──► V = X·W_V (T x d_v) ──────────────────────────────────────────┴──► Salida Z (T x d_v)



\end{lstlisting}

\subsection{4. Atención Multicabeza (\emph{Multi-Head Self-Attention})}

\subsubsection{4.1 Motivación}

Una sola cabeza de atención solo proyecta en un espacio de subespacio único y puede capturar un único tipo de relación (p. ej., sujeto-verbo). El lenguaje requiere capturar múltiples relaciones simultáneamente (sintácticas, semánticas, modificadores).

\subsubsection{4.2 Ecuaciones de Multi-Head Attention}

Se entrenan \(h\) conjuntos independientes de matrices de proyección \(\{W_Q^{(i)}, W_K^{(i)}, W_V^{(i)}\}\):

\[\text{head}_i = \text{Attention}\left(Q^{(i)}, K^{(i)}, V^{(i)}\right) = \text{Attention}\left(X W_Q^{(i)}, \, X W_K^{(i)}, \, X W_V^{(i)}\right)\]

Las salidas de todas las cabezas se concatenan y se proyectan linealmente con \(W_O \in \mathbb{R}^{h d_v \times d_{\text{model}}}\) para mantener las dimensiones de entrada:

\[\text{MultiHead}(X) = \text{Concat}(\text{head}_1, \text{head}_2, \dots, \text{head}_h) W_O\]

\subsubsection{4.3 Dimensiones de Tensores en Práctica (Batching)}

\begin{itemize}

\item Entrada \(X \in \mathbb{R}^{B \times T \times d_{\text{model}}}\) (\(B\): Batch size, \(T\): Longitud de secuencia).

\end{itemize}

\begin{itemize}

\item Proyección y cambio de forma para procesamiento paralelo:

\end{itemize}

\[Q, K, V \rightarrow \mathbb{R}^{B \times h \times T \times d_k}\]

\begin{itemize}

\item Salida de la atención por cabeza: \(\mathbb{R}^{B \times h \times T \times d_v}\).

\end{itemize}

\begin{itemize}

\item Concatenación y proyección final: \(\mathbb{R}^{B \times T \times d_{\text{model}}}\).

\end{itemize}

\subsection{5. Arquitectura del Transformer: Encoder}

El \emph{Encoder} mapea la secuencia de entrada en una representación contextualizada manteniendo la dimensión del eje temporal \(T\) y la dimensión del embedding \(d_{\text{model}}\). Consta de un apilamiento de \(N\) capas idénticas (\(N=6\) en la arquitectura original).

\begin{lstlisting}

Entrada (Tokens) ──► [ Input Embedding ] ──► (+) ──► [ Stack de N Capas Encoder ] ──► Representación Contextualizada

                                             ▲

                                    [ Positional Encoding ]



\end{lstlisting}

\subsubsection{5.1 Codificación Posicional (\emph{Positional Encoding})}

Dado que la autoatención es invariante a permutaciones (trata las palabras como un conjunto no ordenado), se debe añadir explícitamente información de orden.

\begin{itemize}

\item \textbf{Incorporación:} Se suma directamente al embedding de entrada en el primer nivel: \(z_i = x_i + p_{\text{pos}}\).

\end{itemize}

\begin{itemize}

\item \textbf{Sinusoidal Encoding (Original Transformer):}

\end{itemize}

\[PE_{(\text{pos}, 2i)} = \sin\left(\frac{\text{pos}}{10000^{2i / d_{\text{model}}}}\right)\]

\[PE_{(\text{pos}, 2i+1)} = \cos\left(\frac{\text{pos}}{10000^{2i / d_{\text{model}}}}\right)\]

\begin{itemize}

\item \textbf{Learned Positional Embeddings (BERT, GPT):} Tabla de parámetros entrenables \(P \in \mathbb{R}^{T_{\text{max}} \times d_{\text{model}}}\).

\end{itemize}

\subsubsection{5.2 Conexiones Residuales y Normalización por Capa (Add \& Norm)}

Se aplican alrededor de cada subcapa (atención y red feedforward):

\[\text{Output} = \text{LayerNorm}(x + \text{Sublayer}(x))\]

\begin{itemize}

\item \textbf{Residual Connections:} Añaden la entrada de la capa a la salida (\(x + \text{Sublayer}(x)\)), facilitando el flujo del gradiente sin desvanecimiento ni explosión.

\end{itemize}

\begin{itemize}

\item \textbf{Layer Normalization:} Normaliza a través de las características \(d_{\text{model}}\) independientemente para cada token.

\end{itemize}

\begin{itemize}

\item \textbf{Variantes de Orden:}

\item \emph{Post-LN (Original):} \(\text{LN}(x + \text{Sublayer}(x))\).

\end{itemize}

\begin{itemize}

\item \emph{Pre-LN (Moderno - GPT/BERT):} \(x + \text{Sublayer}(\text{LN}(x))\), mejora la estabilidad en redes muy profundas.

\end{itemize}

\subsubsection{5.3 Red Feedforward (FFN)}

Se aplica de forma independiente a cada token (\emph{position-wise}):

\[\text{FFN}(x) = \text{Activation}(x W_1 + b_1) W_2 + b_2\]

\begin{itemize}

\item Expande la dimensión intermedia a \(d_{ff}\) (típicamente \(d_{ff} = 4 \times d_{\text{model}}\)).

\end{itemize}

\begin{itemize}

\item \textbf{Funciones de Activación:}

\item \emph{ReLU:} \(\max(0, x)\) (Transformer original).

\end{itemize}

\begin{itemize}

\item \emph{GELU:} \(x \cdot \Phi(x)\) (Modelos modernos como BERT o GPT).

\end{itemize}

\subsection{6. Arquitectura del Transformer: Decoder}

El \emph{Decoder} genera la secuencia objetivo token a token en modo autorregresivo.

\begin{lstlisting}

Entrada desplazada ──► [ Output Embedding + PE ] ──► [ Masked Self-Attention ] ──► [ Cross-Attention ] ──► [ FFN ] ──► [ Linear + Softmax ]



\end{lstlisting}

\subsubsection{6.1 Masked Multi-Head Self-Attention (Atención Causal)}

Para evitar que el modelo acceda a tokens futuros durante el entrenamiento, se aplica una \textbf{máscara causal} sobre la matriz de puntuaciones previa al Softmax:

\[\text{scores}_{\text{masked}} = \text{scores} + M, \quad \text{donde } M_{ij} = \begin{cases} 0 & \text{si } i \ge j \\ -\infty & \text{si } i < j \end{cases}\]

Esto asegura que en la posición \(t\), la predicción solo dependa de los tokens anteriores \(y_{<t}\).

\subsubsection{6.2 Cross-Attention (Atención Cruzada Encoder-Decoder)}

Conecta la representación del \emph{Encoder} con el \emph{Decoder}:

\begin{itemize}

\item \textbf{Query (\(Q\)):} Proviene del \emph{Decoder} (\(X_{\text{dec}} W_Q\)).

\end{itemize}

\begin{itemize}

\item \textbf{Key (\(K\)) y Value (\(V\)):} Provienen de la salida final del \emph{Encoder} (\(H_{\text{enc}} W_K\), \(H_{\text{enc}} W_V\)).

\end{itemize}

\[\text{CrossAttention}(Q, K, V) = \text{Softmax}\left( \frac{Q K^\top}{\sqrt{d_k}} \right) V\]

\subsubsection{6.3 Capas Finales y Proyección de Salida}

\begin{enumerate}

\item La salida oculta del decodificador \(H_{\text{dec}} \in \mathbb{R}^{B \times T_y \times d_{\text{model}}}\) pasa por una capa lineal que proyecta al tamaño del vocabulario \(V\):

\end{enumerate}

\[\text{logits} = H_{\text{dec}} W_{\text{vocab}}^\top + b\]

\begin{enumerate}

\item Se aplica \textbf{Softmax} para obtener la distribución de probabilidad \(P(y_t \mid y_{<t}, x)\).

\end{enumerate}

\subsection{7. Consejos de Entrenamiento, Inferencia y Arquitecturas}

\subsubsection{7.1 Estrategias de Entrenamiento}

\begin{itemize}

\item \textbf{Subword Tokenization:} Uso de algoritmos como BPE (Byte-Pair Encoding) o WordPiece para evitar la explosión de vocabulario y manejar palabras fuera del diccionario (OOV).

\end{itemize}

\begin{itemize}

\item \textbf{Padding Masking:} Dado que los batches se rellenan con \texttt{<PAD>} para igualar longitudes, se aplica un enmascaramiento con \(-\infty\) para evitar atender a posiciones de relleno.

\end{itemize}

\begin{itemize}

\item \textbf{Warm-up Scheduling:} Incremento gradual de la tasa de aprendizaje durante los primeros pasos seguido de un decaimiento (p. ej., coseno).

\end{itemize}

\subsubsection{7.2 Estrategias de Decodificación en Inferencia}

\begin{itemize}

\item \textbf{Greedy Decoding:} Selecciona \(y_t = \arg\max_v P(v)\). Rápido, pero puede caer en subóptimos locales.

\end{itemize}

\begin{itemize}

\item \textbf{Beam Search:} Mantiene las \(k\) mejores hipótesis parciales acumulando sus probabilidades logarítmicas.

\end{itemize}

\begin{itemize}

\item \textbf{Temperature Scaling:} Ajusta los logits dividiéndolos por \(T\) antes del Softmax para controlar la creatividad/aleatoriedad:

\end{itemize}

\begin{itemize}

\item \(T < 1\): Distribución más picuda (más determinista).

\end{itemize}

\begin{itemize}

\item \(T > 1\): Distribución más plana (más diversa).

\end{itemize}

\subsubsection{7.3 Evolución de las Arquitecturas Transformer}

\begin{itemize}

\item \textbf{Encoder-Only (p. ej., BERT):} Utiliza atención bidireccional. Ideal para tareas de comprensión, clasificación de textos y extracción de entidades.

\end{itemize}

\begin{itemize}

\item \textbf{Decoder-Only (p. ej., GPT-3, LLaMA, Mistral):} Utiliza atención causal. Orientado a la generación de texto autorregresiva.

\end{itemize}

\begin{itemize}

\item \textbf{Encoder-Decoder (p. ej., T5, BART):} Conserva la arquitectura original seq2seq para traducción y resumen.

\end{itemize}

\subsection{8. Práctica en Código (PyTorch)}

A continuación se incluye una implementación ejecutable de un módulo de \textbf{Multi-Head Self-Attention} nativo con enmascaramiento causal opcional:

\begin{lstlisting}[language=Python]

import torch

import torch.nn as nn

import torch.nn.functional as F

import math



class MultiHeadSelfAttention(nn.Module):

    def __init__(self, d_model=512, num_heads=8):

        super(MultiHeadSelfAttention, self).__init__()

        assert d_model % num_heads == 0, "d_model debe ser divisible por num_heads"

        

        self.d_model = d_model

        self.num_heads = num_heads

        self.d_k = d_model // num_heads

        

        # Proyecciones lineales para Q, K, V

        self.W_q = nn.Linear(d_model, d_model, bias=False)

        self.W_k = nn.Linear(d_model, d_model, bias=False)

        self.W_v = nn.Linear(d_model, d_model, bias=False)

        

        # Proyección de salida

        self.W_o = nn.Linear(d_model, d_model, bias=False)



    def forward(self, x, mask=None):

        batch_size, seq_len, _ = x.size()

        

        # 1. Proyección lineal -> (B, T, h, d_k) -> Transposición a (B, h, T, d_k)

        Q = self.W_q(x).view(batch_size, seq_len, self.num_heads, self.d_k).transpose(1, 2)

        K = self.W_k(x).view(batch_size, seq_len, self.num_heads, self.d_k).transpose(1, 2)

        V = self.W_v(x).view(batch_size, seq_len, self.num_heads, self.d_k).transpose(1, 2)

        

        # 2. Scaled Dot-Product Attention

        scores = torch.matmul(Q, K.transpose(-2, -1)) / math.sqrt(self.d_k)

        

        if mask is not None:

            scores = scores.masked_fill(mask == 0, -1e9)

            

        attn_weights = F.softmax(scores, dim=-1)

        attn_output = torch.matmul(attn_weights, V) # (B, h, T, d_k)

        

        # 3. Concatenación de cabezas y proyección de salida

        attn_output = attn_output.transpose(1, 2).contiguous().view(batch_size, seq_len, self.d_model)

        output = self.W_o(attn_output)

        

        return output



# Ejemplo de prueba rápida:

if __name__ == "__main__":

    B, T, d_model = 2, 5, 64

    x = torch.randn(B, T, d_model)

    

    # Máscara Causal (Triangular Inferior)

    causal_mask = torch.tril(torch.ones(T, T)).unsqueeze(0).unsqueeze(0)

    

    mha = MultiHeadSelfAttention(d_model=d_model, num_heads=4)

    out = mha(x, mask=causal_mask)

    print(f"Dimensión del tensor de entrada: {x.shape}")

    print(f"Dimensión del tensor procesado: {out.shape}")



\end{lstlisting}

Dime si necesitas hacer alguna prueba de autoevaluación o repasar algún detalle específico de las fuentes del máster.
\end{document}
